\chapter{结论与未来工作}
\label{cha:conclusion}


本文在 Non-Web（WASI）场景下，在统一的编译与运行配置条件下，对 MicroBenchmark 与 PolyBench 两类程序进行了系统性对比实验，并基于 LLVM IR 提取可解释的静态特征，分析了程序结构特征与性能比值（$ratio_{jit/native}$、$ratio_{aot/native}$）及离散标签之间的关联关系。

综合实验结果，本文得到以下主要结论：

（1）\text{程序结构与性能差异具有显著关联。} 宿主交互密集、动态内存分配频繁以及含有大量不可预测分支的程序结构，与 Wasm 相对于 native 的显著性能劣势具有较高共现性，表明执行边界开销、内存管理路径及不规则控制流是影响性能差距的重要因素。

（2）\text{计算强度与访存模式对性能差异具有调制作用。} 在 PolyBench 数据集中，算术主导型的大规模数值核更易出现较大的 Wasm/native 比值；统计分析表明，$compute\_density$ 与 $compute\_mem\_ratio$ 与性能比值呈正相关，而 $ls\_ratio$ 呈负相关，说明计算强度越高，Wasm 的相对劣势越明显，而访存占比提高则有助于缩小性能差距。

（3）\text{相对性能受热点规模与结构形态共同影响。} 在规模较小且结构规整的热点程序中，Wasm 的执行效率可以接近甚至略优于 native，表明性能差异不仅取决于程序类别，还受到热点规模及后端优化空间的共同影响。

（4）\text{静态特征具备一定预测能力但存在局限。} 在有限样本条件下，基于静态特征构建的二分类模型能够在一定程度上区分 \path{native-better} 与 \path{nonnative-better}，验证了静态结构信息在性能判别中的有效性；然而，其预测能力尚不足以支持高置信度的自动化决策，更适合作为辅助分析与初步筛查工具。

基于上述工作，未来研究可从以下几个方面进一步拓展：

（1）\text{扩展数据规模与多样性。} 引入更多真实应用程序、不同输入规模及多样化 benchmark，以检验当前结论的外部有效性与稳健性。

（2）\text{解耦运行时与编译器因素。} 在多种 Wasm 运行时与不同编译优化等级下重复实验，分别分析字节码表示、运行时实现及编译器优化对性能差异的影响。

（3）\text{融合动态信息提升解释能力。} 结合动态剖析数据与硬件性能计数器（如缓存未命中率、分支预测失效率等），与静态特征进行联合建模，以提升对性能差异来源的解释精度。

（4）\text{引入更贴近实际需求的建模方法。} 探索概率校准与代价敏感学习方法，使模型输出能够与实际可容忍性能损失等工程约束相结合。

（5）\text{构建多目标决策框架。} 将性能之外的因素（如安全隔离、可移植性及维护成本）纳入统一分析框架，从而支持更加全面的技术选型与部署决策。

总体而言，本文从程序结构角度出发，对 WebAssembly 与 native 性能差异进行了系统性分析，并验证了静态特征在解释与预测性能行为方面的潜在价值。相关结论为理解 Wasm 的性能边界及其在工程实践中的适用性提供了参考依据。
